% !PRINTED 628
\weberpage{628}
%<<CONT
zu der Basis $\xi\eta_i$ vermittelt wird, können wir nach der Bezeichnung (9), (10) des vorigen Paragraphen so darstellen:
\weq{4}{(\xi\eta_i) = Y(\eta_i),}
und demnach ergibt sich durch Übergang zu der Basis $\eta_i'$:
\[
(\xi\eta_i') = LYL^{-1}(\eta_i').
\]

Die Funktionen $b_1, b_2, \ldots, b_n$ sind also durch die Funktion $\xi$ vollständig bestimmt.

Die Funktion
\weq{5}{(-1)^n b_n =
\begin{vmatrix}
y_{1,1}, & y_{1,2}, & \ldots, & y_{1,n} \\
y_{2,1}, & y_{2,2}, & \ldots, & y_{2,n} \\
\cdot & \cdot & \cdot & \cdot \\
y_{n,1}, & y_{n,2}, & \ldots, & y_{n,n}
\end{vmatrix}}
heißt die \gesperrt{Norm der Funktion} $\xi$ und wird mit $N(\xi)$ bezeichnet. Über sie gelten folgende Sätze:

1. \gesperrt{Wenn $\xi$ nicht identisch Null ist, so ist $N(\xi)$ von Null verschieden.}

Denn wenn die Determinante des Systems (1) verschwindet, so lassen sich rationale Funktionen $y_1, y_2, \ldots, y_n$, die nicht alle verschwinden, so bestimmen, daß
\[
\xi(y_1\eta_1 + y_2\eta_2 + \cdots + y_n\eta_n) = 0
\]
wird, und dies fordert, da $\eta_1, \eta_2, \ldots, \eta_n$ eine Basis ist, $\xi = 0$.

2. \gesperrt{Ist $\xi$ eine rationale Funktion von $z$, so ist ihre Norm die $n$te Potenz dieser Funktion.}

Denn ist $\xi$ rational, so reduzieren sich die Gleichungen (1) auf die Identitäten $\xi\eta_i = \xi\eta_i$. Es verschwinden also in der Determinante (5) alle Glieder mit Ausnahme der Diagonalglieder, und diese werden alle gleich $\xi$. Wir drücken diesen Satz durch die Formel aus:
\[
N(a) = a^n.
\]

3. \gesperrt{Sind $\xi, \xi'$ zwei Funktionen in $\Omega$, so gilt der Satz}
\weq{6}{N(\xi\xi') = N(\xi)N(\xi').}

Denn ist nach (4):
\[
(\xi\eta_i) = Y(\eta_i), \quad (\xi'\eta_i) = Y'(\eta_i),
\]
so ist
\[
(\xi\xi'\eta_i) = YY'(\eta_i),
\]
und daraus folgt die Formel (6), weil die Determinante einer zusammengesetzten linearen Substitution, hier $YY'$, gleich dem Produkt der Determinante der Komponenten ist.
% [sic] printed "Produkt der Determinante der Komponenten" (singular "Determinante"), the sense requires the plural "Determinanten": the determinant of a composed substitution is the product of the determinants of the components
